\section{Evaluation Plan}
\label{sec:evaluation}

\paragraph{Reference infrastructures.}
Evaluating capture and replication requires ground truth that production systems rarely provide, and we will construct a
suite of reference networks for which the topology, the configuration, and the set of planted
vulnerabilities are fully known to us. These networks will be built from
widely deployed open-source components, including our anchor products and the
databases, brokers, identity providers, and orchestrators that surround them,
and they will be described entirely through infrastructure-as-code, so we can instantiate
them reproducibly, vary their scale, and generate the
ground-truth knowledge graph automatically. We have developed substantial
expertise in this domain through our work on the Global Accessible Test
Infrastructure (GATE) as part of the ACTION NSF AI
Institute~\cite{action}, in which we have created topologies that represent
smart cities, chemical plants, water systems, and manufacturing plants. We
will seed each reference network with a mixture of known vulnerabilities that
are drawn from public advisories, including the nginx defect classes of
Section~\ref{sec:landscape}, and with deliberately introduced
misconfigurations. In addition, we will design a subset of the networks to
require genuine multi-step chains, so success cannot be achieved by finding a
single isolated flaw. These reference networks will also serve as
the challenges for the twinCTF competition described in
Section~\ref{sec:broader}, and they will be released as a public benchmark.

\paragraph{Metrics and targets.}
Targets are stated for the reference networks, where ground truth exists.

{\footnotesize\setlength{\tabcolsep}{4pt}\renewcommand{\arraystretch}{0.93}
\begin{longtable}{@{}p{0.55\linewidth}p{0.42\linewidth}@{}}
\toprule
\textbf{Metric} & \textbf{Target}\\
\midrule
\multicolumn{2}{@{}l}{\emph{Thrust 1 --- capture}}\\
Hosts and services correctly identified & $\geq 95\%$ of ground truth\\
Service dependencies correctly identified & $\geq 90\%$\\
Model coverage estimate vs.\ measured & within 10 points\\
Expert input per 100 hosts & $\leq 4$ h, declining across cycles\\
Production impact during observation & no measurable latency or error change\\
\addlinespace[2pt]
\multicolumn{2}{@{}l}{\emph{Thrust 2 --- replication}}\\
Attack-path recall (paths preserved) & $\geq 90\%$\\
Precision (twin paths real in original) & $\geq 75\%$\\
Resource reduction vs.\ full replica & $\geq 10\times$ at equal recall\\
\addlinespace[2pt]
\multicolumn{2}{@{}l}{\emph{Thrust 3 --- analysis}}\\
Seeded single-component defects found & $\geq 85\%$\\
Seeded multi-step chains found & $\geq 70\%$\\
Findings failing to reproduce on original & $\leq 20\%$\\
Run-to-run stability, identical inputs & $\geq 0.8$ Jaccard overlap\\
Cost per assessment cycle (100 hosts) & reported; reduced $\geq 2\times$\\
\textbf{Confirmed defects reported upstream} & $\geq 10$ across nginx, ZMap, deps\\
\bottomrule
\end{longtable}}

We regard the upstream-report target as the most important line in this table,
since a fix that lands upstream reaches every deployment of the affected
component.

\paragraph{Baselines and ablations.}
Where established baselines exist, we will compare against them. For
component-level analysis, we will compare the defects found by \sysname
against those found by running the underlying fuzzing, static-analysis, and
symbolic-execution tools in isolation, in order to quantify the value added
by the agentic orchestration and by the context of the knowledge base. For
capture, we will compare against an unmodified ZMap~\cite{zmap} and against
inventory tools such as osquery~\cite{osquery} and
cartography~\cite{cartography}, measuring both coverage and production
impact, and, for configuration properties, against network configuration
verification~\cite{batfish}. For multi-step analysis, we will compare against
our own ChainReactor~\cite{chainreactor} on the single-host problems for
which it was designed, and against attack-graph generation~\cite{mulval} on
networked problems. Throughout, we will report the human effort required at
each stage, since a central claim of the project is that \sysname shifts
experts from manual assessment to the verification of high-confidence
findings.

\paragraph{Validation with existing users.}
Reference networks establish correctness against ground truth; whether the
results are useful can only be determined on real deployments. Working with the collaborators whose letters accompany this
proposal, and with industry and critical-infrastructure partners, we will
evaluate \sysname on operational systems in the later phase of the project.
The ZMap extensions of Section~\ref{sec:thrust1} will be validated by
operators who run them on their own networks and who report coverage and
impact, and the nginx configuration checks of Section~\ref{sec:thrust3} will
be validated against real deployed configurations contributed by
partners. Subject to written authorization, as described in
Section~\ref{sec:buildtest}, we will apply the full framework to the internal
research networks at UCSB and ASU and to partner systems. In these settings,
the emphasis of the evaluation shifts to the non-disruptiveness of the
observation phase, to the accuracy of the continuously maintained model as
the infrastructure changes, and to whether a finding changed what the
administrators who own the systems actually did.

\paragraph{Robustness to the evolution of language models.}
The design of \sysname separates the components that are independent of any
model from the judgments that are delegated to one. The knowledge-base schema
with its provenance and confidence annotations, the equivalence-class
abstraction, the fidelity measure, and the separation of passive observation
from aggressive analysis are all specified without reference to how the agents
are implemented, and they constitute the research contribution of this
project. The agents supply the judgments that those specifications leave open,
such as interpreting an unfamiliar configuration dialect, reconciling a manual
that contradicts a runtime observation, and composing exploitation primitives
into a chain. The two activities that dominate the quality of an assessment,
namely the capture of an infrastructure's configuration and the discovery of
the flaws it contains, both fall on the judgment side, so an improvement in
model capability yields a more complete knowledge base and longer chains with
no change to the schema, to the observation planner, or to the planning
formulation. Because the agents interact only through typed interfaces, models
are swappable, and we will evaluate open-weight models alongside frontier ones
so that the framework does not depend on any single provider. The reference
networks carry ground truth, which lets us hold the benchmark fixed and vary
the model. We will run that substitution across models of differing capability
at each annual review, and it tests the separation as much as it tests the
models: if some substitution requires a change to the schema, to the
observation planner, or to the planning formulation, then the boundary between
structure and judgment has been drawn in the wrong place, and we will report
where it belongs. The same gains in capability accrue to the adversary, and
the substitution experiment measures how much of each gain the defender
collects.
